% part: intuitionistic-logic
% chapter: propositions-as-types
% section: types

\documentclass[../../../include/open-logic-section]{subfiles}

\begin{document}

\olfileid[tr]{int}{pty}{typ}

\olsection{Tipler}

$(MN)$ gibi bir ispat terimi, hangi ispatın ispat terimi olduğunu tek
başına biricik biçimde belirlemez. Bunun nedeni, serbest değişkenlerin, yani
aksiyomların ispat terimlerinin karşılık gelen aksiyomlardaki !!{formula}s{}i
belirtmemesidir. Ancak bunları bir $\Gamma$ bağlamında belirtirsek ispat
\emph{biricik} biçimde belirlenir.

\begin{defn}
Bir ispat terimi~$N$ için \emph{bağlam}, $N$'nin her serbest değişkeni için
tam olarak bir $x : !A$ çifti içeren bir $\Gamma$ kümesidir.
\end{defn}

Tanımın, $\Gamma$'nın \emph{yalnız} $N$'nin serbest değişkenleri için
$x:!A$ içermesini gerektirmediğine dikkat edin. Dolayısıyla örneğin
$\Gamma = \{x:!A, y:!B, z:!C\}$, $\pair{x}{y}$ için bir bağlamdır.

\begin{defn}\ollabel{defn:type}
  $\Gamma$'nın $N$ için bir bağlam olduğunu varsayın. $N$'nin
  ($\Gamma$'ya göre) \emph{tipi} tümevarımsal olarak şöyle tanımlanır:
  \begin{enumerate}
  \item $x$'in tipi ancak ve ancak $x:!A \in \Gamma$ ise $!A$'dır.
  \item $N$'nin $x:!A, \Gamma$'ya göre tipi~$!B$ ise
  $\lambd[\typeof{x}{!A}][N]$'nin tipi $!A \lif !B$'dir.
  \item $N$'nin tipi $!A \lif !B$ ve $M$'nin tipi~$!A$ ise $(NM)$'nin
  tipi~$!B$'dir.
  \item $N$'nin tipi $!A$ ve $M$'nin tipi~$!B$ ise $\pair{N}{M}$'nin
  tipi $!A \land !B$'dir.
  \item $N$'nin tipi~$!A_1 \land !A_2$ ise $\proj{i}{N}$'nin tipi~$!A_i$'dir.
  \item $N$'nin tipi $!A$ ise $i=1$ olduğunda $\inj[!B]{i}{N}$'nin tipi
  $!A \lor !B$; $i=2$ olduğunda $!B \lor !A$'dır.
  \item $M$'nin tipi $!A \lor !B$, $N_1$'in $x:!A, \Gamma$'ya göre tipi
  $!C$ ve $N_2$'nin $y:!B, \Gamma$'ya göre tipi~$!C$ ise
  $\dcase{M}{x}{N_1}{y}{N_2}$'nin tipi~$!C$'dir.
  \item $N$'nin tipi~$\lfalse$ ise $\abort{!A}{N}$'nin tipi~$!A$'dır.
  \end{enumerate}
\end{defn}

\begin{prop}
Her ispat terimi~$N$, kendisi için bir bağlam~$\Gamma$'ya göre en fazla
bir tipe sahiptir; başka bir deyişle, iyi tiplenmiş her ispat teriminin tipi
biriciktir.
\end{prop}

\begin{proof}
  $N$ üzerinde tümevarımla.
\end{proof}

$N$'nin $\Gamma$'ya göre tipi~$!A$ ise $\Gamma \Sequent N : !A$ yazarız.
\olref{defn:type} içindeki koşullar tanıdık gelmelidir: küçük bir farkla tam
olarak \Log{tN2i}'nin kurallarıdır; bağlam~$\Gamma$ baştan sona aynıdır.
Tanımı, \olref{tab:tN3ip} içinde verilen kurallarla bir \Log{tN3i} hesabı
olarak kodlayabiliriz.

\subfile{rules-tN3}

İspat terimleri, lambda kalkülüsünün bir sürümü olan \emph{çarpımlar,
toplamlar ve boş tip içeren tipli lambda kalkülüsü}ndeki terimlerdir.
Tiplenmemiş lambda kalkülüsü yalnız değişkenlerden, \emph{lambda
soyutlamaları}~$\lambd[x][N]$ ve uygulama $(MN)$ ile kurulan terimlere
sahiptir; tipli lambda soyutlaması $\lambd[\typeof{x}{!A}][N]$ içindeki
~$!A$ açıklamasına sahip değildir. Lambda kalkülüsü Turing-tam bir hesaplama
modelidir (yani her kısmi hesaplanabilir fonksiyon onda bir terimle temsil
edilebilir). İspat terimleri ve tipleme kurallarının verdiği lambda
kalkülüsü sürümü sınırlı bir hesaplama modelidir; ancak temel düşünceler
aynıdır. Tip atamasıyla sınırlandırılmıştır: yalnız \emph{iyi tiplenmiş}
terimler, yani (bir $\Gamma$ bağlamına göre) tipi bulunan ispat terimleri
yasal terimlerdir. Örneğin $\lambd[x][(xx)]$, tiplenmemiş lambda
kalkülüsünde yasal bir terimdir; ancak $\lambd[\typeof{x}{!A}][(xx)]$ hiçbir
$!A$ için iyi tiplenmiş değildir. Çünkü tipleme kuralları, $(xx)$'in bir
tipe sahip olabilmesi için ilk~$x$'in tipinin $!A \lif !B$, ikinci $x$'in
tipinin~$!A$ olmasını gerektirir. $!A$, $!A \lif !B$'nin bir öz
alt-!!{formula}{}ü olduğundan bu olanaksızdır.

Tipleri sezgisel olarak, o tipteki bir terimin adlandırdığı şeyin tipi gibi
düşünebiliriz. Dolayısıyla $!A \land !B$ tipindeki bir terim, ilk bileşeni
$!A$, ikinci bileşeni~$!B$ tipinde olan bir çifti, yani $!A \times !B$
çarpımının bir !!a{element}{}ını seçer. $!A \lif !B$ tipindeki bir terim,
$!A$ tipindeki şeylerden $!B$ tipindeki şeylere giden bir fonksiyondur.
$!A \lor !B$ tipindeki bir terim, $!A$ tipindeki ya da $!B$ tipindeki bir
şey ile hangisi olduğuna ilişkin bilgiden oluşur. Bu, $!A$ ile $!B$
kümelerinin $\{\tuple{i,x} : i = 1 \text{ veya } 2, i = 1 \text{ ise }
x \in !A, i=2 \text{ ise } x \in !B\}$ olarak tanımlanan ayrık birleşimine
benzer. $\lfalse$ tipi boş kümedir; yani bir terim bu tipe sahipse hiçbir
şeyi adlandıramaz. Doğal tümdengelim (\Log{tN3i} dâhil) tutarlı olduğundan,
açık varsayımlar bulunmadıkça~$\lfalse$'ın hiçbir !!a{proof}{}ı yoktur;
dolayısıyla tipi~$\lfalse$ olan kapalı bir ispat terimi (serbest değişkensiz
bir terim) yoktur.

Tipli lambda kalkülüsünün işleçleri, uygulandıkları terimlerin belirli
şeyleri adlandırması koşuluyla sezgisel olarak doğru tipte bir şeyi adlandıran
terimler üretir. Örneğin ``$N_1 : !A$ ve $N_2 : !B$ ise
$\pair{N_1}{N_2} : !A \land !B$'', $N_1$ tipi~$!A$ olan bir şeyi ve $N_2$
tipi~$!B$ olan bir şeyi adlandırıyorsa $\pair{N_1}{N_2}$'nin tipi
$!A \land !B$ olan bir şeyi adlandırdığını söyler.

\end{document}
